\chapter{Introdução}
\label{cap:01}

O desenvolvimento de jogos digitais representa uma área estratégica na interseção entre tecnologia, arte e entretenimento. Dentro dessa ampla indústria, destacam-se os \textit{cozy games}, um subgênero que prioriza experiências relaxantes, visuais acolhedores e mecânicas acessíveis \cite{waszkiewicz2020}. Ao contrário de jogos voltados para ação ou competição intensa, os \textit{cozy games} propõem uma interação tranquila, centrada em rotinas simples, ambientações reconfortantes e vínculos afetivos com o ambiente virtual.

A escolha por esse tema partiu do interesse em compreender como elementos técnicos e narrativos podem ser combinados para criar experiências imersivas voltadas ao conforto emocional. Observa-se uma carência na literatura referente ao processo de desenvolvimento desse gênero de jogo, especialmente no uso da \textit{Godot Engine}, uma plataforma de código aberto que tem ganhado destaque entre desenvolvedores independentes \cite{GodotDocs2024}. Este trabalho pretende, assim, contribuir para esse campo pouco explorado.

 Com base nisso, o objetivo principal deste trabalho é criar um jogo 2D do gênero \textit{cozy} usando a \textit{engine Godot}, buscando uma experiência ao mesmo tempo imersiva e relaxante. Para alcançar isso, os objetivos específicos incluem definir o estilo visual e a narrativa, desenvolver mecânicas que promovam uma sensação de conforto, preparar o ambiente de desenvolvimento, selecionar e adaptar os recursos visuais e sonoros, além de implementar as funcionalidades fundamentais do jogo.

A metodologia adotada baseia-se na abordagem prática de estudo de caso, com etapas iterativas de desenvolvimento. São utilizados referenciais técnicos como \citeauthorandyear{Bradfield2018} e a documentação oficial da \citeauthorandyear{GodotDocs2024}, além de autores que discutem a relação entre regras, narrativa e experiência estética em jogos, como \citeauthorandyear{Juul2005} e \citeauthorandyear{Suits1978}. O escopo do trabalho limita-se à criação de um protótipo funcional, priorizando coerência temática e estética em detrimento de complexidade técnica.


\section{Objetivos}

\subsection{Objetivo Geral}

Desenvolver o jogo 2D "Cat's Village", um jogo \textit{cozy game} no estilo \textit{Top Down} utilizando a \textit{engine Godot}, com o intuito de aplicar e demonstrar conhecimentos sobre desenvolvimento de jogos.

\subsection{Objetivos Específicos}
\begin{itemize}
	\item Definir o estilo visual e narrativo do jogo;
	\item Projetar jogabilidade e interações relaxantes;
	\item Preparar o ambiente de desenvolvimento na \textit{Godot};
	\item Criar ou adquirir os \textit{assets} visuais e sonoros;
	\item Programar as mecânicas e sistemas do jogo.
	
\end{itemize}